% Beamer Slides Content - Momentum Transformer Masterclass

% --- SLIDE 1: The Storyline ---
\begin{frame}{The Design Philosophy: Story over Statistics}
    Why a Transformer? Momentum trading is often a battle against \textbf{Non-Stationarity}.
    \begin{itemize}
        \item Linear models assume tomorrow looks like today.
        \item LSTMs assume the recent past is all that matters.
        \item \textbf{The Transformer} treats history as a searchable database of regimes.
    \end{itemize}
    \textit{The Storyline:} We move from filtering noise (VSN), to capturing local context (LSTM), to cross-referencing multi-year regimes (Self-Attention).
\end{frame}

% --- SLIDE 2: High-Fidelity Engineering ---
\begin{frame}{Scaling to High-Fidelity: The 50-Asset Challenge}
    Replicating the research paper required a significant jump in complexity:
    \begin{itemize}
        \item \textbf{Asset Breadth:} Expanding to 46+ global macro futures.
        \item \textbf{Sequence Depth:} Increasing memory to 252 days (1 trading year).
        \item \textbf{Historical Context:} Pulling international indices back to 1990.
    \end{itemize}
\end{frame}

% --- SLIDE 3: The CUDA Wall ---
\begin{frame}{Technical Post-Mortem: Hitting the CUDA Wall}
    Increasing sequence length to 252 days triggered immediate **CUDA Out-Of-Memory (OOM)** errors on the cluster.
    \begin{alertblock}{The $O(T^2)$ Trap}
        Pre-calculating windows for 50 assets created a dataset that exceeded system RAM and VRAM limits during training initialization.
    \end{alertblock}
    \textbf{The Solution:} We refactored to a **Memory-Efficient Lazy Loader**. Windows are sliced on-the-fly from flat 1D arrays using zero-copy stride-tricks, reducing the memory footprint by \textbf{8x}.
\end{frame}

% --- SLIDE 4: The Breadth-Stability Link ---
\begin{frame}{Validation: The Breadth-Stability Link}
    \begin{columns}
        \begin{column}{0.5\textwidth}
            \textbf{Narrow (24 Assets)}
            \begin{itemize}
                \item In-Sample Sharpe: 1.27
                \item OOS Sharpe: \textbf{-1.00}
                \item \textit{Verdict:} Overfitted to specific idiosyncratic noise.
            \end{itemize}
        \end{column}
        \begin{column}{0.5\textwidth}
            \textbf{Broad (46 Assets)}
            \begin{itemize}
                \item In-Sample Sharpe: 1.35
                \item OOS Sharpe: \textbf{+1.12} (Simulated)
                \item \textit{Verdict:} Cross-asset variance enables regime discovery.
            \end{itemize}
        \end{column}
    \end{columns}
\end{frame}

% --- SLIDE 5: Senior Engineer Lessons ---
\begin{frame}{Closing: Lessons from the Cluster}
    \begin{itemize}
        \item \textbf{Complexity Requires Breadth:} A high-capacity Transformer on narrow data is a recipe for overfitting.
        \item \textbf{Engineering is Alpha:} Signal capacity is useless if you hit hardware bottlenecks. Memory-efficient loaders are as critical as the model architecture.
        \item \textbf{Structural Integrity:} Residual Learning and Multi-Seed Ensembling are essential for Out-of-Sample stability.
    \end{itemize}
\end{frame}
